\subsection{Hệ thống dữ liệu}
Dữ liệu của dự án đến từ hai bộ dữ liệu công khai của Olist trên Kaggle, được mô phỏng thành hai hệ thống nguồn độc lập trên PostgreSQL. Hệ thống thứ nhất là hệ thống bán hàng (OLTP), gồm 9 bảng với 1.550.922 dòng, ghi nhận đơn hàng, sản phẩm đã bán, thanh toán, đánh giá, khách hàng, người bán, danh mục sản phẩm và toạ độ địa lý theo mã bưu chính. Hệ thống thứ hai là hệ thống CRM/Marketing, gồm 2 bảng với 8.842 dòng: 8.000 khách tiềm năng (MQL) và 842 thương vụ đã chốt. Hai hệ thống chỉ gặp nhau qua mã người bán \code{seller\_\allowbreak{}id}; vì là hai hệ thống khác nhau nên cơ sở dữ liệu không đặt khoá ngoại giữa chúng và việc liên kết được thực hiện ở kho dữ liệu thông qua chiều \code{dim\_\allowbreak{}seller}. Tổng dung lượng 11 file CSV là 121,18 MB (1.559.764 dòng), riêng bảng toạ độ địa lý chiếm 58,44 MB và hơn một triệu dòng. Mỗi hệ thống nằm trong một schema riêng (\code{oltp}, \code{crm}), giữ nguyên tên cột gốc, kể cả hai cột sai chính tả \code{product\_\allowbreak{}name\_\allowbreak{}lenght} và \code{product\_\allowbreak{}description\_\allowbreak{}lenght}, vì đó là đặc điểm của hệ thống nguồn và chỉ được sửa ở tầng ODS. Dữ liệu được nạp nguyên trạng, không làm sạch. Nhóm tạo hai database: \code{olist\_\allowbreak{}full} (262,01 MB) chứa toàn bộ dữ liệu để khám phá, và \code{olist\_\allowbreak{}dw} (185,14 MB) chỉ chứa các bản ghi đến hết 31/12/2017 (1.273.620 dòng) để làm điểm xuất phát cho nạp tăng dần. Số dòng nạp vào \code{olist\_\allowbreak{}full} khớp hoàn toàn với số dòng trong file CSV và không có dòng nào bị loại.

\begin{table}[htbp]
\centering\small
\caption{Hai hệ thống nguồn: số bảng ghi, số cột và dung lượng file CSV}
\label{tab:src-tables}
\begin{tabular}{llrrr}
\toprule
\textbf{Bảng nguồn} & \textbf{File CSV} & \textbf{Số dòng} & \textbf{Số cột} & \textbf{MB} \\
\midrule
oltp.product\_category\_name\_translation & product\_category\_name\_translation.csv & 71 & 2 & 0,00 \\
oltp.products & olist\_products\_dataset.csv & 32.951 & 9 & 2,27 \\
oltp.sellers & olist\_sellers\_dataset.csv & 3.095 & 4 & 0,17 \\
oltp.geolocation & olist\_geolocation\_dataset.csv & 1.000.163 & 6 & 58,44 \\
oltp.customers & olist\_customers\_dataset.csv & 99.441 & 5 & 8,62 \\
oltp.orders & olist\_orders\_dataset.csv & 99.441 & 8 & 16,84 \\
oltp.order\_items & olist\_order\_items\_dataset.csv & 112.650 & 7 & 14,72 \\
oltp.order\_payments & olist\_order\_payments\_dataset.csv & 103.886 & 5 & 5,51 \\
oltp.order\_reviews & olist\_order\_reviews\_dataset.csv & 99.224 & 7 & 13,78 \\
crm.mql & olist\_marketing\_qualified\_leads\_dataset.csv & 8.000 & 4 & 0,67 \\
crm.closed\_deals & olist\_closed\_deals\_dataset.csv & 842 & 14 & 0,16 \\
Tổng (11 bảng) &  & 1.559.764 &  & 121,18 \\
\bottomrule
\end{tabular}
\end{table}

Hai hệ thống nguồn được mô phỏng thành hai schema riêng (\code{oltp} và \code{crm}) trong PostgreSQL; sau khi nạp đủ, kho có số dòng từng lớp như Bảng \ref{tab:layers}.
\begin{table}[htbp]
\centering\small
\caption{Số bảng và số dòng của từng lớp trong kho dữ liệu (trạng thái cuối)}
\label{tab:layers}
\begin{tabular}{llrr}
\toprule
\textbf{Lớp} & \textbf{Schema} & \textbf{Số bảng} & \textbf{Số dòng} \\
\midrule
Nguồn bán hàng & oltp & 9 & 1.550.922 \\
Nguồn marketing & crm & 2 & 8.842 \\
Staging & staging & 11 & 1.559.764 \\
ODS & ods & 18 & 678.183 \\
Kho (dimension + fact) & dw & 16 & 614.223 \\
Giám sát ETL & etl\_log & 6 & 2.701 \\
Data mart (view và materialized view) & mart\_* & 17 view & -- \\
\bottomrule
\end{tabular}
\end{table}

\paragraph{Tần suất và khối lượng cập nhật.} Dữ liệu Kaggle là ảnh chụp tĩnh nên nhóm mô phỏng cập nhật theo tháng: nạp ban đầu đến 31/12/2017 (45.430 đơn) rồi 11 đợt hằng tháng (01--11/2018), tổng 286.144 bản ghi mới (xem mục 3.4.3).
\todo{Nhóm nêu rõ tần suất cập nhật giả định của hệ thống thực (ví dụ đơn hàng cập nhật liên tục, kho cập nhật hằng ngày) và dung lượng mỗi lần cập nhật.}
